\documentclass[10pt, aspectratio=169]{beamer}

% --- PAKIETY ---
\usepackage[utf8]{inputenc}
\usepackage[T1]{fontenc}
\usepackage[polish]{babel}
\usepackage{amsmath, amsfonts, amssymb}
\usepackage{graphicx}
\usepackage{hyperref}
\usepackage{multimedia} % Do osadzania wideo
\usepackage{tikz}
\usetikzlibrary{angles, quotes, calc, 3d, arrows.meta, positioning}

% --- WYGLĄD PREZENTACJI ---
\usetheme{UWr_white}
% \usecolortheme{seahorse} % Usunięto ze względu na własny motyw
% \setbeamertemplate{navigation symbols}{} % Znajduje się w motywie UWr

% --- DANE TYTUŁOWE ---
\title[Matematyka i Drony]{Matematyczne podstawy modelowania trajektorii lotu dronów}
\subtitle{Od równań na tablicy do autonomicznej misji}
\author{Maciej Tadej}
\institute{Wydział Matematyki i Informatyki \\ Uniwersytet Wrocławski}
\date{Dolnośląski Kongres Młodego Odkrywcy}

\begin{document}

\begin{frame}
    \titlepage
\end{frame}

\begin{frame}{Plan prezentacji}
    \tableofcontents
\end{frame}

% --- SEKCJE PREZENTACJI ---
\section{Wstęp}

\begin{frame}{Kiedy myślimy o dronach...}
    \begin{columns}
        \column{0.5\textwidth}
        Zazwyczaj wyobrażamy sobie operatora z padem w dłoni. 
        \vspace{0.5cm}
        
        Jednak w nowoczesnym przemyśle, ratownictwie i logistyce sterowanie ręczne to przeszłość.
        
        \column{0.5\textwidth}
        \begin{center}
        \begin{tikzpicture}
            \node[inner sep=0pt] (drone) at (0,0) {\includegraphics[width=0.8\textwidth]{example-image-a}}; % Miejsce na obrazek drona/operatora
            \node[below=0.2cm of drone] {Od sterowania do autonomii};
        \end{tikzpicture}
        \end{center}
    \end{columns}
\end{frame}

\begin{frame}{Schemat quadcoptera -- jak to działa?}
    Aby maszyna mogła latać, potrzebuje czterech wirników, które generują ciąg (siłę nośną).
    \vspace{0.5cm}
    
    \begin{center}
    \begin{tikzpicture}[scale=0.8]
        % Ramiona drona
        \draw[line width=3pt, gray!80] (-2,-2) -- (2,2);
        \draw[line width=3pt, gray!80] (-2,2) -- (2,-2);
        
        % Korpus
        \filldraw[blue!70!black] (0,0) circle (0.8cm) node[white, font=\bfseries] {DRON};
        
        % Wirniki (z zaznaczeniem kierunku obrotu)
        % Górny lewy - CW
        \filldraw[gray!30, draw=black, thick] (-2,2) ellipse (1cm and 0.3cm);
        \draw[->, thick, red] (-2, 2.5) arc (90:-90:0.5cm and 0.2cm);
        
        % Górny prawy - CCW
        \filldraw[gray!30, draw=black, thick] (2,2) ellipse (1cm and 0.3cm);
        \draw[->, thick, green!70!black] (2, 2.5) arc (90:270:0.5cm and 0.2cm);
        
        % Dolny lewy - CCW
        \filldraw[gray!30, draw=black, thick] (-2,-2) ellipse (1cm and 0.3cm);
        \draw[->, thick, green!70!black] (-2, -1.5) arc (90:270:0.5cm and 0.2cm);
        
        % Dolny prawy - CW
        \filldraw[gray!30, draw=black, thick] (2,-2) ellipse (1cm and 0.3cm);
        \draw[->, thick, red] (2, -1.5) arc (90:-90:0.5cm and 0.2cm);
        
        % Oznaczenia
        \node[red] at (-2, 3) {Kierunek 1};
        \node[green!70!black] at (2, 3) {Kierunek 2};
    \end{tikzpicture}
    \end{center}
\end{frame}

\begin{frame}{Czym jest autonomia?}
    Autonomia w robotyce oznacza, że maszyna sama podejmuje decyzje i wie:
    \begin{itemize}
        \item Gdzie się znajduje?
        \item Gdzie ma polecieć?
        \item Jak tam dotrzeć, omijając przeszkody?
    \end{itemize}
    \vspace{0.5cm}
    \textbf{Odpowiedzią na te pytania jest matematyka.} Pojęcia z lekcji matematyki i fizyki stają się potężnymi narzędziami, z których korzystają współczesne technologie.
\end{frame}

\begin{frame}{Inżynieria lotu -- Jak drony komunikują się ze światem?}
    Za autonomią stoi skomplikowany, ale bardzo uporządkowany proces komunikacji. 
    Drony porozumiewają się ze stacją naziemną (GCS) za pomocą protokołu \textbf{MAVLink}.
    
    \vspace{0.4cm}
    \begin{center}
    \begin{tikzpicture}[scale=0.9, transform shape,
        box/.style={draw, rounded corners, fill=blue!10, text width=3cm, align=center, minimum height=1.2cm},
        arrow/.style={->, thick, >=stealth}]
        
        % Komponenty
        \node[box, fill=green!20] (gcs) at (0, 0) {\textbf{Stacja Bazowa (GCS)} \\ \footnotesize QGroundControl \\ Mission Planner};
        \node[box, fill=orange!20] (mavlink) at (5, 0) {\textbf{Protokół MAVLink} \\ \footnotesize Telemetria, Komendy};
        \node[box, fill=gray!20] (drone) at (10, 0) {\textbf{Kontroler Lotu} \\ \footnotesize ArduPilot (SITL) \\ Pixhawk};
        
        % Strzałki komunikacji
        \draw[arrow, bend left=15] (gcs.east) to node[above, font=\footnotesize]{SET\_MODE (GUIDED)} (mavlink.west);
        \draw[arrow, bend left=15] (mavlink.east) to node[above, font=\footnotesize]{ARM (Uzbrojenie)} (drone.west);
        
        \draw[arrow, bend left=15] (drone.south west) to node[below, font=\footnotesize]{GPS, Attitude (Pitch/Roll)} (mavlink.south east);
        \draw[arrow, bend left=15] (mavlink.south west) to node[below, font=\footnotesize]{HEARTBEAT (Status)} (gcs.south east);
        
        % Proces Pre-Arm
        \node[draw, dashed, fill=yellow!10, text width=4cm, align=left] at (10, -2.5) {
            \textbf{Pre-Arm Checks:}\\
            1. GPS Lock (3D Fix)\\
            2. IMU Calibration\\
            3. Barometer OK\\
            4. Bateria > 20\%
        };
        \draw[->, dashed, thick] (10, -1.2) -- (10, -0.6);
        
    \end{tikzpicture}
    \end{center}
\end{frame}

\section{Podstawy ruchu}

\begin{frame}{Globalne współrzędne -- Kula (Układ Sferyczny)}
    Zanim dron poleci, musi wiedzieć gdzie jest na Ziemi. Najprostszym modelem jest \textbf{kula}. Pozycję określamy za pomocą długości i szerokości geograficznej oraz promienia.
    
    \vspace{0.2cm}
    \begin{center}
    \begin{tikzpicture}[scale=1.5]
        % Sfera
        \draw[thick] (0,0) circle (1.5cm);
        \draw[dashed, gray] (0,0) ellipse (1.5cm and 0.4cm); % Równik tył
        \draw[gray] (-1.5,0) arc (180:360:1.5cm and 0.4cm);  % Równik przód
        
        % Oś Z
        \draw[->, thick] (0,-1.5) -- (0,2) node[anchor=south]{$Z$ (Biegun Północny)};
        
        % Punkt
        \filldraw[blue] (0.8, 0.8) circle (1.5pt) node[anchor=south west]{Dron $(r, \theta, \phi)$};
        \draw[->, thick, blue] (0,0) -- (0.8, 0.8) node[midway, above left]{$r$};
        
        % Rzuty
        \draw[dashed] (0.8, 0.8) -- (0.8, 0.15);
        \draw[dashed] (0,0) -- (0.8, 0.15);
        
        % Kąty
        \draw[red, thick] (0,0.5) arc (90:45:0.5cm) node[midway, right]{$\theta$};
    \end{tikzpicture}
    \end{center}
\end{frame}

\begin{frame}{Czy Ziemia to idealna kula? -- Układ Elipsoidalny}
    Ziemia jest spłaszczona na biegunach. W systemach GPS (np. standardzie WGS84) dron korzysta z \textbf{elipsoidy obrotowej}. 
    
    \vspace{0.2cm}
    \begin{center}
    \begin{tikzpicture}[scale=1.5]
        % Elipsoida
        \draw[thick] (0,0) ellipse (1.8cm and 1.2cm);
        \draw[dashed, gray] (0,0) ellipse (1.8cm and 0.3cm); 
        \draw[gray] (-1.8,0) arc (180:360:1.8cm and 0.3cm); 
        
        % Osie
        \draw[->, thick] (0,-1.3) -- (0,1.8) node[anchor=south]{Oś obrotu};
        \draw[->, thick] (0,0) -- (2.2,0) node[anchor=west]{Równik ($a$)};
        
        % Punkt i normalna
        \filldraw[blue] (1.1, 0.8) circle (1.5pt) node[anchor=south west]{Dron};
        \draw[->, red, thick] (1.1, 0.8) -- (1.5, 1.2) node[anchor=south]{Normalna do powierzchni};
        \draw[dashed] (1.1, 0.8) -- (0.5, 0);
        
        \node[font=\footnotesize] at (0, -1.6) {WGS84: $a \approx 6378$ km (równik), $b \approx 6356$ km (biegun)};
    \end{tikzpicture}
    \end{center}
\end{frame}

\begin{frame}{Gdy jesteśmy blisko celu -- Układ lokalny NED}
    Globalne współrzędne (GPS) mają dziwne stopnie i radiany. Dla wygody obliczeń, wokół drona tworzy się płaski układ lokalny \textbf{NED} (North-East-Down).
    
    \vspace{0.2cm}
    \begin{center}
    \begin{tikzpicture}[scale=1.5, x={(0.866cm, -0.5cm)}, y={(0.866cm, 0.5cm)}, z={(0cm, -1cm)}]
        % Osie NED (Z skierowane w dół!)
        \draw[->, thick, red] (0,0,0) -- (2,0,0) node[anchor=west]{$N$ (Północ / $X$)};
        \draw[->, thick, green!60!black] (0,0,0) -- (0,2,0) node[anchor=south]{$E$ (Wschód / $Y$)};
        \draw[->, thick, blue] (0,0,0) -- (0,0,2) node[anchor=north]{$D$ (W dół / $Z$)};
        
        % Dron
        \filldraw[black] (0,0,0) circle (2pt) node[anchor=south, xshift=-0.5cm]{Start (0,0,0)};
        \filldraw[orange] (1.5, 1.5, -1) circle (2pt) node[anchor=south west]{Cel $(x, y, -h)$};
        
        \draw[dashed, gray] (0,0,0) -- (1.5, 1.5, 0) -- (1.5, 1.5, -1);
    \end{tikzpicture}
    \end{center}
\end{frame}

\begin{frame}{Start drona -- pierwszy kod}
    Korzystając z lokalnego układu NED, polecamy maszynie polecieć 10 metrów na Północ i 5 metrów w Górę ($D = -5$).
    
    \vspace{0.5cm}
    \textbf{Demonstracja wideo:} Start i lot po prostej trajektorii.
    
    \begin{center}
    \begin{tikzpicture}
        \draw[fill=gray!20] (0,0) rectangle (6,3);
        \node at (3,1.5) {MIEJSCE NA WIDEO: Start i prosta trajektoria};
    \end{tikzpicture}
    \end{center}
\end{frame}

\section{Fizyka i matematyka lotu}

\begin{frame}{Siły i Wektory - Wiatr jako wyzwanie}
    W idealnym świecie drony latałyby po idealnie prostych liniach. W rzeczywistości przeszkadza im wiatr. Wiatr to z punktu widzenia matematyki nic innego jak \textbf{wektor}.
    
    \vspace{0.5cm}
    \begin{center}
    \begin{tikzpicture}[scale=1.5]
        \draw[thick, ->, blue] (0,0) -- (2,1) node[midway, above left]{Prędkość drona $\vec{V_d}$};
        \draw[thick, ->, red] (2,1) -- (3,0) node[midway, above right]{Wiatr $\vec{W}$};
        \draw[thick, ->, dashed, green!50!black] (0,0) -- (3,0) node[midway, below]{Rzeczywisty tor lotu $\vec{V_{rz}}$};
    \end{tikzpicture}
    \end{center}
    \vspace{0.2cm}
    $\vec{V_{rz}} = \vec{V_d} + \vec{W}$. By utrzymać kurs, dron musi korygować swój lot.
\end{frame}

\begin{frame}{Trajektorie - krzywe parametryczne}
    Do płynnego lotu, np. nagrywania kamerą obiektu w centrum, drony poruszają się po okręgach lub innych kształtach. Do ich opisania używamy równań krzywych.
    
    \vspace{0.3cm}
    \textbf{Przykład - Lot po krzywej w kształcie serca:}
    \begin{columns}
        \begin{column}{0.5\textwidth}
            \small
            Używamy równań parametrycznych (zmienna $t$ od $0$ do $2\pi$):
            \begin{align*}
                x(t) &= 16 \sin^3(t) \\
                y(t) &= 13 \cos(t) - 5 \cos(2t) \\
                     &\quad - 2 \cos(3t) - \cos(4t)
            \end{align*}
            Komputer pokładowy dzieli tę krzywą na mniejsze odcinki i wysyła je jako kolejne współrzędne (waypointy) w formacie geograficznym (LAT/LON).
        \end{column}
        \begin{column}{0.5\textwidth}
            \begin{center}
            \begin{tikzpicture}[scale=0.15]
                \draw[gray!30, thin, step=5] (-20,-20) grid (20,15);
                \draw[->] (-20,0) -- (20,0) node[right]{$x$};
                \draw[->] (0,-20) -- (0,15) node[above]{$y$};
                
                \draw[thick, red, domain=0:2*pi, samples=100, variable=\t] 
                    plot ({16*sin(\t r)^3}, {13*cos(\t r) - 5*cos(2*\t r) - 2*cos(3*\t r) - cos(4*\t r)});
                
                \node[blue, font=\tiny] at (16, 5) {$\bullet$ Dron};
            \end{tikzpicture}
            \end{center}
        \end{column}
    \end{columns}
\end{frame}

\begin{frame}{Dyskretyzacja trajektorii}
    Aby dron mógł wykonać lot, ciągła funkcja matematyczna musi zostać zamieniona na \textbf{punkty orientacyjne (waypointy)}.
    Dzielimy czas $t$ na równe interwały $\Delta t$.
    
    \vspace{0.3cm}
    \begin{columns}
        \begin{column}{0.5\textwidth}
            \small
            Pomiędzy punktami dron porusza się po liniach prostych (najkrótszą drogą). 
            Im gęściej próbkujemy (mniejsze $\Delta t$), tym ruch jest płynniejszy, ale wymaga przesłania do kontrolera lotu większej liczby danych.
            
            \vspace{0.3cm}
            Tutaj krzywą podzielono na 20 równych odcinków czasowych:
            \begin{align*}
                t_i &= i \cdot \Delta t \\
                \Delta t &= \frac{2\pi}{20}
            \end{align*}
        \end{column}
        \begin{column}{0.5\textwidth}
            \begin{center}
            \begin{tikzpicture}[scale=0.15]
                \draw[gray!30, thin, step=5] (-20,-20) grid (20,15);
                \draw[->] (-20,0) -- (20,0) node[right]{$x$};
                \draw[->] (0,-20) -- (0,15) node[above]{$y$};
                
                % Cienka, przezroczysta idealna krzywa w tle
                \draw[gray!50, thin, domain=0:2*pi, samples=100, variable=\t] 
                    plot ({16*sin(\t r)^3}, {13*cos(\t r) - 5*cos(2*\t r) - 2*cos(3*\t r) - cos(4*\t r)});
                    
                % Segmenty linii - zdyskretyzowana trajektoria
                \draw[thick, blue, domain=0:2*pi, samples=21, variable=\t] 
                    plot ({16*sin(\t r)^3}, {13*cos(\t r) - 5*cos(2*\t r) - 2*cos(3*\t r) - cos(4*\t r)});
                
                % Kropki (Waypointy) - wyliczane na podstawie tego samego wzoru
                \foreach \i in {0, 1, ..., 20} {
                    \pgfmathsetmacro{\t}{\i * pi / 10}
                    \pgfmathsetmacro{\x}{16*sin(\t r)^3}
                    \pgfmathsetmacro{\y}{13*cos(\t r) - 5*cos(2*\t r) - 2*cos(3*\t r) - cos(4*\t r)}
                    \filldraw[red] (\x, \y) circle (0.6);
                }
            \end{tikzpicture}
            \end{center}
        \end{column}
    \end{columns}
\end{frame}

\section{Misje autonomiczne}

\begin{frame}{Skanowanie obszaru i optymalizacja tras}
    Ratownik szukający zaginionej osoby nie steruje dronem ręcznie nad całym lasem. Zadaje mu misję skanowania obszaru.
    \vspace{0.2cm}
    Jak wyznaczyć optymalną trasę z punktu $A$ do punktu $B$, zahaczając o inne ważne punkty? Z pomocą przychodzi teoria grafów i Problem Komiwojażera (TSP - Traveling Salesperson Problem).
    
    \begin{center}
    \begin{tikzpicture}[node distance=1.5cm, every node/.style={circle, draw, minimum size=0.5cm}]
        \node (1) at (0,0) {A};
        \node (2) [above right of=1] {B};
        \node (3) [below right of=2] {C};
        \node (4) [below right of=1] {D};
        
        \path[->, thick, blue]
            (1) edge (2)
            (2) edge (3)
            (3) edge (4)
            (4) edge (1);
    \end{tikzpicture}
    \end{center}
\end{frame}

\begin{frame}{Zaawansowane misje ratunkowe 3D}
    W rzeczywistości teren rzadko bywa płaski. Drony muszą uwzględniać wysokość (3D). 
    \vspace{0.5cm}
    
    \textbf{Demonstracja wideo:} Skanowanie obszaru poligonu, przeszukiwanie objętości cylindra i pościg za celem ruchomym.
    
    \begin{center}
    \begin{tikzpicture}
        \draw[fill=gray!20] (0,0) rectangle (7,3);
        \node at (3.5,1.5) {MIEJSCE NA WIDEO: Algorytmy szukania 3D};
    \end{tikzpicture}
    \end{center}
\end{frame}

\section{Roje dronów}

\begin{frame}{Jak ptaki i owady stają się inspiracją?}
    Zamiast jednego, potężnego drona, z zadaniem lepiej potrafi poradzić sobie rój kilkunastu (lub kilkuset!) małych jednostek. Ale jak uniknąć zderzeń w powietrzu? 
    \vspace{0.3cm}
    
    Opisujemy to matematycznym \textbf{modelem Cuckera-Smale'a}, który opiera się na prostych zasadach:
    \begin{itemize}
        \item \textbf{Separacja}: nie wpadaj na sąsiada.
        \item \textbf{Kohezja}: trzymaj się blisko grupy.
        \item \textbf{Aliniacja (Dopasowanie)}: leć w tym samym kierunku, co reszta.
    \end{itemize}
\end{frame}

\begin{frame}{Rój w środowisku ArduPilot}
    Samo oprogramowanie nie wie początkowo, że ma do czynienia z wieloma maszynami. Zestawiamy architekturę opartą o rozproszoną komunikację między systemami.
    
    \vspace{0.3cm}
    Drony wymieniają się swoim wektorem położenia i prędkości.
    
    \begin{columns}
        \begin{column}{0.5\textwidth}
            \small
            \textbf{Architektura Komunikacji} \\
            Każdy dron wysyła swoją pozycję (GPS) i prędkość przez protokół MAVLink. 
            Algorytm roju odbiera te pakiety, liczy odległości do sąsiadów na podstawie wzorów, i wysyła korekty wektora prędkości.
        \end{column}
        \begin{column}{0.5\textwidth}
            \begin{center}
            \begin{tikzpicture}[scale=0.8, every node/.style={transform shape}]
                \node (d1) at (0,2) [circle, draw, fill=blue!20] {Dron 1};
                \node (d2) at (4,2) [circle, draw, fill=blue!20] {Dron 2};
                \node (d3) at (2,0) [circle, draw, fill=blue!20] {Dron 3};
                
                \draw[<->, dashed, thick] (d1) -- (d2) node[midway, above, font=\tiny] {MAVLink (Pozycja)};
                \draw[<->, dashed, thick] (d2) -- (d3);
                \draw[<->, dashed, thick] (d3) -- (d1);
                
                \node (algo) at (2,4) [rectangle, draw, fill=green!20] {Algorytm Roju (Cucker-Smale)};
                \draw[->, thick, red] (algo) -- (d1) node[midway, left, font=\tiny] {Korekta $\vec{V}$};
                \draw[->, thick, red] (algo) -- (d2);
                \draw[->, thick, red] (algo) -- (d3);
            \end{tikzpicture}
            \end{center}
        \end{column}
    \end{columns}
\end{frame}

\begin{frame}{Rozdzielenie roju i kontrola odległości}
    W obliczu przeszkody (np. wzniesienie, budynek), ławica musi podjąć decyzję o rozdzieleniu i późniejszym połączeniu.
    
    \vspace{0.3cm}
    Rozwiązując w każdym ułamku sekundy układ równań różniczkowych, drony nieustannie utrzymują zaprogramowany margines odległości (bezpieczeństwa) od innych maszyn.
    
    \begin{columns}
        \begin{column}{0.5\textwidth}
            \small
            \textbf{Omijanie Przeszkód (Boids / Siły Odrzucające)} \\
            Gdy na drodze pojawia się przeszkoda, algorytm generuje wokół niej "pole siłowe" (odpychające). 
            Drony reagują na tę wirtualną siłę rozdzielając się na dwie grupy. 
            Gdy przeszkoda znika, siły kohezji (przyciągania) ponownie łączą rój w jedną ławicę.
        \end{column}
        \begin{column}{0.5\textwidth}
            \begin{center}
            \begin{tikzpicture}[scale=0.8]
                % Przeszkoda
                \node (obs) at (2.5,2.5) [circle, draw, fill=red!40, minimum size=1.2cm] {O};
                \draw[red, dashed, thick] (2.5,2.5) circle (1.2cm);
                
                % Drony omijajace
                \node (d1) at (0,1) [circle, fill=blue, inner sep=1.5pt] {};
                \draw[->, thick, blue] (d1) to[out=45, in=180] (1.5, 3.5) to[out=0, in=135] (4, 2.5);
                
                \node (d2) at (0,4) [circle, fill=blue, inner sep=1.5pt] {};
                \draw[->, thick, blue] (d2) to[out=-45, in=180] (1.5, 1.5) to[out=0, in=-135] (4, 2.5);
                
                \node (d3) at (4,2.5) [circle, fill=green, inner sep=1.5pt] {};
                \draw[->, thick, green] (d3) -- (5,2.5) node[right, font=\tiny] {Połączony rój};
                
            \end{tikzpicture}
            \end{center}
        \end{column}
    \end{columns}
\end{frame}

\section{Podsumowanie}

\begin{frame}{Podsumowanie}
    \begin{itemize}
        \item Programowanie dronów wymaga patrzenia na świat przez pryzmat liczb, geometrii i fizyki.
        \item Skomplikowane zachowania, takie jak lot roju ptaków, można sprowadzić do układu równań.
        \item Matematyka, której uczycie się w szkole, nie jest tylko zbiorem teoretycznych zagadek. To narzędzie, które pozwala technologii samodzielnie "myśleć" i działać.
    \end{itemize}
    
    \vspace{0.5cm}
    \begin{center}
        \Large \textbf{Dziękuję za uwagę!} \\
        \normalsize Czas na Wasze pytania.
    \end{center}
\end{frame}


\end{document}
